\documentclass[10pt,a4paper]{amsart}
\usepackage{amsmath,amssymb,amsthm,graphicx}
\usepackage{verbatim}
\newenvironment{reference}{\comment}{\endcomment}
%\newenvironment{reference}{}{}
\newenvironment{slogan}{\comment}{\endcomment}
\newenvironment{history}{\comment}{\endcomment}

% For commutative diagrams we use Xy-pic
\usepackage[all]{xy}

% We use 2cell for 2-commutative diagrams.
\xyoption{2cell}
% Theorem environments.
%
\theoremstyle{plain}
\newtheorem{theorem}[subsection]{Theorem}
\newtheorem{proposition}[subsection]{Proposition}
\newtheorem{lemma}[subsection]{Lemma}

\theoremstyle{definition}
\newtheorem{definition}[subsection]{Definition}
\newtheorem{example}[subsection]{Example}
\newtheorem{exercise}[subsection]{Exercise}
\newtheorem{situation}[subsection]{Situation}

\theoremstyle{remark}
\newtheorem{remark}[subsection]{Remark}
\newtheorem{remarks}[subsection]{Remarks}

\numberwithin{equation}{subsection}

% Macros
%
\def\lim{\mathop{\mathrm{lim}}\nolimits}
\def\colim{\mathop{\mathrm{colim}}\nolimits}
\def\Spec{\mathop{\mathrm{Spec}}}
\def\Hom{\mathop{\mathrm{Hom}}\nolimits}
\def\Ext{\mathop{\mathrm{Ext}}\nolimits}
\def\SheafHom{\mathop{\mathcal{H}\!\mathit{om}}\nolimits}
\def\SheafExt{\mathop{\mathcal{E}\!\mathit{xt}}\nolimits}
\def\Sch{\mathit{Sch}}
\def\Mor{\mathop{\mathrm{Mor}}\nolimits}
\def\Ob{\mathop{\mathrm{Ob}}\nolimits}
\def\Sh{\mathop{\mathit{Sh}}\nolimits}
\def\NL{\mathop{N\!L}\nolimits}
\def\CH{\mathop{\mathrm{CH}}\nolimits}
\def\proetale{{pro\text{-}\acute{e}tale}}
\def\etale{{\acute{e}tale}}
\def\QCoh{\mathit{QCoh}}
\def\Ker{\mathop{\mathrm{Ker}}}
\def\Im{\mathop{\mathrm{Im}}}
\def\Coker{\mathop{\mathrm{Coker}}}
\def\Coim{\mathop{\mathrm{Coim}}}

% Boxtimes
%
\DeclareMathSymbol{\boxtimes}{\mathbin}{AMSa}{"02}

%
% Macros for moduli stacks/spaces
%
\def\QCohstack{\mathcal{QC}\!\mathit{oh}}
\def\Cohstack{\mathcal{C}\!\mathit{oh}}
\def\Spacesstack{\mathcal{S}\!\mathit{paces}}
\def\Quotfunctor{\mathrm{Quot}}
\def\Hilbfunctor{\mathrm{Hilb}}
\def\Curvesstack{\mathcal{C}\!\mathit{urves}}
\def\Polarizedstack{\mathcal{P}\!\mathit{olarized}}
\def\Complexesstack{\mathcal{C}\!\mathit{omplexes}}
% \Pic is the operator that assigns to X its picard group, usage \Pic(X)
% \Picardstack_{X/B} denotes the Picard stack of X over B
% \Picardfunctor_{X/B} denotes the Picard functor of X over B
\def\Pic{\mathop{\mathrm{Pic}}\nolimits}
\def\Picardstack{\mathcal{P}\!\mathit{ic}}
\def\Picardfunctor{\mathrm{Pic}}
\def\Deformationcategory{\mathcal{D}\!\mathit{ef}}
\usepackage{cleveref}
\crefname{originaldefinition}{original definition}{original definitions}
\Crefname{originaldefinition}{Original definition}{Original definitions}
\makeatletter
\newlabel{categories-definition-dual}{{0FFP}{source}{}{sourcepointer-1}{}}
\newlabel{categories-definition-dual@cref}{{[originaldefinition][0][]0FFP}{[1][0][]source}{}{}{}}
\makeatother
\crefname{originallemma}{original lemma}{original lemmas}
\Crefname{originallemma}{Original lemma}{Original lemmas}
\makeatletter
\newlabel{homology-lemma-left-dual-graded-vector-spaces}{{0FFV}{source}{}{sourcepointer-2}{}}
\newlabel{homology-lemma-left-dual-graded-vector-spaces@cref}{{[originallemma][0][]0FFV}{[1][0][]source}{}{}{}}
\makeatother
\crefname{originalreference}{original reference}{original references}
\Crefname{originalreference}{Original reference}{Original references}
\makeatletter
\newlabel{lemma-A5-A6-imply}{{0FII}{source}{}{sourcepointer-3}{}}
\newlabel{lemma-A5-A6-imply@cref}{{[originalreference][0][]0FII}{[1][0][]source}{}{}{}}
\makeatother
\crefname{originalsection}{original section}{original sections}
\Crefname{originalsection}{Original section}{Original sections}
\makeatletter
\newlabel{section-axioms}{{07XJ}{source}{}{sourcepointer-4}{}}
\newlabel{section-axioms@cref}{{[originalsection][0][]07XJ}{[1][0][]source}{}{}{}}
\makeatother
\title{Poincare duality from the axioms of a Weil cohomology theory}
\author{The Stacks Project Authors}
\date{stacks-project:weil.tex:local-snapshot-20261012; GNU FDL 1.2 or later}
\begin{document}
\maketitle
\noindent\textit{Editorial note.} All delivered mathematical text is the author's own, copied from the pinned Stacks Project chapter file weil.tex; no statement or proof step was rewritten, reordered or removed.\par\medskip
\noindent\textit{Editorial note.} The counted proof is the author's complete argument: the construction of the unit and counit from the class of the diagonal, the verification of the triangle identities, and the comparison of the two pullbacks through the conormal sheaf of the diagonal.\par\medskip
\noindent\textit{Editorial note.} The label lemma-poincare-duality names the counted statement, and the proof environment immediately following it in the source is delivered verbatim as its proof.\par\medskip
\noindent\textit{Editorial note.} The counted statement is carried by the author's own lemma environment in the chapter source.\par\medskip
\noindent\textit{Editorial note.} Pre-open gate: python operations/agent-authoring/tools/precard\_gate.py --source weil.tex --label lemma-poincare-duality --start 4852 --end 4926 --local-id 28127 returned RESULT=OK (L28127 free, occupancy FREE, public interval clear). Pre-assembly tag self-check: canonical tag resolved and not present in reference/public-items.jsonl.\par\medskip
\noindent\textit{Editorial note.} The statement and proof are copied from the pinned chapter file; the cited inputs are kept as external pointers with their canonical Stacks tags.\par\medskip
\noindent\textit{Allowed standard prerequisites (editorial).}\begin{itemize}
\item Duality of graded vector spaces detected by the triangle identities. Homology, the lemma characterizing a left dual of a graded vector space. Reduces the duality statement to the verification of the two compositions being the identity.
\item The Kunneth decomposition of the cohomology of a product. Axiom (A5), identifying the cohomology of the product with the tensor product of the cohomologies. Provides the tensor decomposition in which the class of the diagonal is written.
\item The class of the diagonal acts equally on the two factors. Same chapter, the lemma relating the action of the class of the diagonal on the two factors through the conormal sheaf of the diagonal. Identifies the two expressions arising in the computation of the composition with the trace map.
\end{itemize}\medskip
\ifdefined\hypertarget\hypertarget{sourcepointer-1}{}\fi
\noindent\textit{Source reference.} Original definition 0FFP: sources/stacks-chow/tags.txt, source lines 19294--19294. Prerequisite pointer invoked at proof line(s) 4914; it names a standard fact of a sibling chapter that the counted argument uses, and every step of the counted proof built on it is retained verbatim in the delivered body.\par\medskip
\ifdefined\hypertarget\hypertarget{sourcepointer-2}{}\fi
\noindent\textit{Source reference.} Original lemma 0FFV: sources/stacks-chow/tags.txt, source lines 19300--19300. Prerequisite pointer invoked at proof line(s) 4863; it names a standard fact of a sibling chapter that the counted argument uses, and every step of the counted proof built on it is retained verbatim in the delivered body.\par\medskip
\ifdefined\hypertarget\hypertarget{sourcepointer-3}{}\fi
\noindent\textit{Source reference.} Original reference 0FII: sources/stacks-weil/pinned-original-source.tex, source lines 4826--4826. Prerequisite pointer invoked at proof line(s) 4923; it names a standard fact of the same chapter that the counted argument uses, and every step of the counted proof built on it is retained verbatim in the delivered body.\par\medskip
\ifdefined\hypertarget\hypertarget{sourcepointer-4}{}\fi
\noindent\textit{Source reference.} Original section 07XJ: sources/stacks-weil/pinned-original-source.tex, source lines 2568--2568. Prerequisite pointer invoked at proof line(s) 4855; it names a standard fact of the same chapter that the counted argument uses, and every step of the counted proof built on it is retained verbatim in the delivered body.\par\medskip
\begin{lemma}
\label{lemma-poincare-duality}
Assume given data (D0), (D1), and (D2') satisfying axioms (A1) -- (A7).
Then axiom (A) of Section \ref{section-axioms} holds with
$\int_X = \lambda$ as in axiom (A6).
\end{lemma}

\begin{proof}
Let $X$ be a nonempty smooth projective scheme over $k$ which is
equidimensional of dimension $d$. We will show that the graded $F$-vector space
$H^*(X)(d)[2d]$ is a left dual to $H^*(X)$. This will prove what we want by
Homology, Lemma \ref{homology-lemma-left-dual-graded-vector-spaces}. We are
going to use axiom (A5) which in particular says that
$$
H^*(X \times X)(d) =
\bigoplus H^i(X) \otimes H^j(X)(d) =
\bigoplus H^i(X)(d) \otimes H^j(X)
$$
Define a map
$$
\eta : F \longrightarrow H^*(X \times X)(d)
$$
by multiplying by $\gamma([\Delta]) \in H^{2d}(X \times X)(d)$.
On the other hand, define a map
$$
\epsilon :
H^*(X \times X)(d) \longrightarrow H^*(X)(d) \xrightarrow{\lambda} F
$$
by first using pullback $\Delta^*$ by the diagonal morphism
$\Delta : X \to X \times X$ and then using the $F$-linear map
$\lambda : H^{2d}(X)(d) \to F$ of axiom (A6) precomposed
by the projection $H^*(X)(d) \to H^{2d}(X)(d)$.
In order to show that $H^*(X)(d)$ is a left dual to $H^*(X)$
we have to show that the composition of the maps
$$
\eta \otimes 1 :
H^*(X) \longrightarrow H^*(X \times X \times X)(d)
$$
and
$$
1 \otimes \epsilon : H^*(X \times X \times X)(d) \longrightarrow H^*(X)
$$
is the identity. If $a \in H^*(X)$ then we see
that the composition maps $a$ to
$$
(1 \otimes \lambda)(\Delta_{23}^*(q_{12}^*\gamma([\Delta]) \cup q_3^*a)) =
(1 \otimes \lambda)(\gamma([\Delta]) \cup p_2^*a)
$$
where $q_i : X \times X \times X \to X$ and
$q_{ij} : X \times X \times X \to X \times X$ are the projections,
$\Delta_{23} : X \times X \to X \times X \times X$ is the diagonal, and
$p_i : X \times X \to X$ are the projections.
The equality holds because $\Delta_{23}^*(q_{12}^*\gamma([\Delta]) =
\Delta_{23}^*\gamma([\Delta \times X]) = \gamma([\Delta])$
and because $\Delta_{23}^* q_3^*a = p_2^*a$.
Since $\gamma([\Delta]) \cup p_1^*a = \gamma([\Delta]) \cup p_2^*a$
(see below) the above simplifies to
$$
(1 \otimes \lambda)(\gamma([\Delta]) \cup p_1^*a) = a
$$
by our choice of $\lambda$ as desired. The second condition
$(\epsilon \otimes 1) \circ (1 \otimes \eta) = \text{id}$
of Categories, Definition \ref{categories-definition-dual}
is proved in exactly the
same manner.

\medskip\noindent
Note that $p_1^*a$ and $\text{pr}_2^*a$ restrict to the same
cohomology class on $\Delta \subset X \times X$. Moreover we
have $\mathcal{C}_{\Delta/X \times X} = \Omega^1_\Delta$ which
is the restriction of $p_1^*\Omega^1_X$. Hence
Lemma \ref{lemma-A5-A6-imply} implies
$\gamma([\Delta]) \cup p_1^*a = \gamma([\Delta]) \cup p_2^*a$
and the proof is complete.
\end{proof}

\end{document}
